\section{Acceptance, replay, and trust}\label{sec:acceptance}

\subsection{A separate acceptance transaction}

The search engine is allowed to be heuristic, buggy, or interrupted. The
acceptance path should be small enough to audit independently:
\begin{enumerate}
\item Instantiate all assigned expression and universe metavariables. Reject
unresolved holes, dangling locals, and unresolved elaboration obligations.
Generalize only the universe parameters and locals permitted by the original
request, not newly invented assumptions.
\item Close the program and contract proof over the exact dependent local
context. Preserve local definitions appropriately, or substitute them with
checked capture-avoiding abstraction. Collect generated helper declarations
in dependency order.
\item Reconstruct the original requested type and predicate from the frozen
request, not from a candidate-supplied summary. Check the closed program and
proof using the ordinary kernel-checked declaration-admission path.
\item Audit the program's and proof's transitive assumptions and the generated
auxiliaries. Reject \code{sorryAx}, unapproved axioms, and any bypass of the
selected kernel-checking policy. Check computational-provider restrictions
separately from proof-provider restrictions.
\item Freeze the accepted expressions, declaration delta, request identity,
and audit result together. Only then create a verified result record and a
source presentation.
\end{enumerate}

An \code{inferType} result or successful \code{isDefEq} call is useful during
search but does not replace declaration admission. In particular, a configured
skip-kernel-check option or an unchecked admission API must not be reachable
through the verified-results path. A compiler run containing errors may insert
placeholder terms; the resulting existence of a declaration name is not a
successful verification.

The theorem established by this path is Equation~\eqref{eq:acceptance}, relative
to the original environment and admitted assumptions. It does not certify the
searcher's completeness, ranking quality, or runtime performance. Writing the
searcher in Lean is compatible with its implementation being \code{partial};
its generated programs are separate declarations with their own dependency
and computability checks.

\subsection{Display is a view, not candidate identity}

An accepted \Expr{} is authoritative. Delaborated syntax is a view for people.
A source-export command should replay the emitted source in the same context
and verify that it has the requested type and contract. When exact preservation
of the accepted program is required, compare the replayed expression with the
original using conversion or a checked equality appropriate to the type.
If that comparison cannot be established, report a newly checked source
variant rather than claim it is the identical accepted object.

Names, implicit arguments, notation, coercions, and local instances can change
how the same text elaborates. Export should therefore qualify ambiguous names,
make delicate universe/type/dictionary selections explicit, and retain
required helper definitions. Source-level simplification is another proposal
that needs replay; it cannot borrow the original candidate's proof by matching
the printed function name.

\subsection{Independent validation options}

A separate process using the same Lean kernel checks operational isolation and
prevents stale metaprogramming state from leaking into acceptance. It is not an
independent implementation of the kernel. For higher assurance, an additional
checker such as Lean4Lean can be evaluated for supported artifacts and pinned
versions, with its own documented scope~\cite{lean4lean}. This is an optional
validation route, not a claim that this article has performed it.

A particularly useful regression is to replay the \emph{exact emitted program}
and contract theorem in a fresh module that omits synthesis tactics and test
oracles from the construction environment. That detects accidental dependency
on helper assumptions, not just syntax problems in the final pretty printer.

\subsection{Receipt structure}

A production receipt should include the original request, environment/toolchain
identity, provider and policy selections, accepted declarations, contract type,
proof dependency closure, runtime dependency observations, source-replay
outcome, search bounds, and outcome counts. It should also distinguish candidate
rejections from rule inapplicability and infrastructure failures.

The receipt's opaque accepted-object identifier is tied to its environment
generation. A fingerprint can detect accidental mismatches but cannot mint a
new accepted object. If the contract changes, create a new request and verify
again. If only ranking changes, existing verified candidates may be reranked
without pretending that new search work or a new correctness proof occurred.

\section{Practical implementation and resource behavior}\label{sec:engineering}

\subsection{Explicit cursors instead of implicit infinite search}

Represent every generator as a finite-step cursor returning a proposal,
continuation, completion, or diagnostic. An explicit cursor is easier to pause,
account for, and serialize than a lazily evaluated stream whose next element
can trigger an unbounded search. It also avoids making Haskell-style laziness
a compatibility requirement for the Lean implementation.

Persistent snapshots permit efficient branch sharing, but retained snapshots
can keep large expression graphs alive. Benchmark memory with realistic
coupled goals, not just short proof terms. Use arenas or generation-indexed
storage for ephemeral metadata, compact decision trails, and bounded caches
whose entries are closed templates rather than whole abandoned worker states.
Do not assume a Lean rewrite is faster or slower without measuring the actual
workload.

The main blow-up is usually branching, not the cost of one constructor
application. If $H$ heads are considered and a head creates $a$ unconstrained
arguments, naive independent enumeration multiplies the choices across all
$a$ positions. Result-directed unification, argument dependencies, specifications,
and observation partitions exist to avoid that product as early as possible.
They are more consequential than micro-optimizing a queue implementation.

\subsection{Budget dimensions}

Maintain separate counters for rule expansions, unification work, generated
AST cost, proof tasks, solver calls, candidate evaluations, live branch memory,
and elapsed time. A global wall-clock deadline is still necessary because not
all operations have comparable heartbeat behavior. Hard memory exhaustion is
an infrastructure outcome, not a completed negative search.

Per-rule limits should yield a suspended or unknown branch where possible.
Increasing a proof budget should not require regenerating every earlier term.
Likewise, increasing a program-cost bound should reuse validated indexes and
proof templates without reusing stale branch-local metavariables. Log the
actual resources charged rather than only the user-configured maximum.

\subsection{Cancellation and process isolation}

A worker owns every external solver process it creates, including descendants.
Cancellation closes solver inputs, terminates the owned process tree when
necessary, and waits for cleanup before reusing any channel. Responses carry
request and solver-frame identities so a delayed reply cannot satisfy a newer
obligation.

The ordinary native engine needs no external solver to produce basic results.
Optional solver launch hardening should be documented independently from
synthesis semantics. Reproducing a predecessor's low-level descriptor-bound
launch guarantees might require a small platform binding; it is not a reason
to reintroduce Haskell as the semantic host. No external tool may install an
axiom simply to make its answer accepted.

\subsection{Determinism and reproducibility}

Deterministic mode uses stable provider ordering, deterministic tie-breaking,
fixed seeds for test generation, and a reproducible logical schedule. Actual
wall-clock deadlines can still create variation; record that limitation.
Performance mode may accept the first verified result from parallel workers,
but the receipt must retain which lane and branch produced it.

Requests should pin the Lean toolchain and imported package versions. API
adaptation belongs in \code{LeanBridge}; search policy should not be scattered
with version-specific field accesses. The included prototype intentionally
pins Lean 4.34.0 rather than claim source compatibility with every Lean release.

\subsection{Diagnostics that explain the actual obstacle}

A useful failure report identifies the nearest residual obligations and the
kind of restriction encountered. Examples are: a dependent index equality not
proved; an unresolved instance output parameter; a required constructor absent
from the selected grammar; a recursive schema requiring generalization; or a
candidate satisfying all examples with its universal proof still open.

Also report successful partial mechanisms without upgrading them to full
success. Finding a fold carrier and its step separately is useful diagnostic
information, but not evidence that the original composition was synthesized.
Showing an actual rejected wrong candidate is stronger negative-control
evidence than merely showing that a false-contract query printed no output.

An editor integration can visualize the obligation graph, show the decision
that constrained a hole, and offer an explicit sketch annotation at that point.
The normal batch result should remain machine-readable and concise: full traces
belong in attached artifacts or an opt-in diagnostic channel.
